\section{Analysis}\label{analysis}

\subsection{newwin}\label{TCLnewwin}

Syntax:

{\tt newwin} {\em name}

Creates a new toplevel window for a TCL widget, with Tk identifier
{\em name}. {\em name} must start with a full stop.

\subsection{Reduction Functions}

Syntax:

{\tt max} {\em tcllist}\\
{\tt min} {\em tcllist}\\
{\tt av} {\em tcllist}\\

The maximum, minimum and average of a TCL list.

\subsection{Palette Variable}\label{palette}

This variable is used by both the {\tt display\_stub} and {\tt
  connect\_stub} routines to define a palette of colours for colouring
  the species in the instruments. If the TCL variable {\tt palette} is
  assigned a list of X-windows colours (on many systems, a list of
  such colours is found in {\tt /usr/lib/X11/rgb.txt}),
  then the palette class can be used like an array within C++,
  returning the colour name as a string:
\begin{verbatim}
palette[i]
\end{verbatim}
returns the \verb|i%n|th colour, where {\tt n} is the number of
colours in the palette list.

\subsection{display\_stub}

Syntax:

\begin{verbatim}
void display_stub(iarray& density, iarray& species, iarray& nsp, 
                  int ncells=1, int cell=0);
\end{verbatim}

This function implements the {\tt display}(\S\ref{display}) widget.
The value of each component of density is plotted as a function of
{\tt tstep}, and coloured by the index {\tt species}. Set {\tt
  ncells>1} and {\tt cell} in order to display only a single cell, as
in ecolab. These variables may be ignored if not meaningful. {\tt nsp}
must be set with the cellular structure used by {\tt density}. see
\S\ref{spatial}. For a non-cellular model, it contains just the 1
component, equal the {\tt density.size}.

\subsection{connect\_stub}

Syntax:

\begin{verbatim}
void connect_stub(iarray& density, sparse_mat& interaction, iarray& nsp, 
                  int ncells=1)
\end{verbatim}

This function implements the {\tt connect\_plot}(\S\ref{connect})
widget. The {\tt row} and {\tt col} connections are plotted, and
sorted and coloured according to connected subsets. {\tt density} is
only used to indicate extinctions (where {\tt density[i]==0}, the row
and column is coloured in wheat. Set {\tt ncells>1} and {\tt cell} in
order to display only a single cell, as in ecolab. These variables may
be ignored if not meaningful. {\tt nsp} must be set with the cellular
structure used by {\tt density} and {\tt interaction}. see
\S\ref{spatial} For a non-cellular model, it contains just the 1
component, equal the {\tt density.size}.

